\section{MoE 为什么流行：更多参数，不一定更多 FLOPs}

本节解释 MoE 的核心动机：同样每 token FLOPs 下增加参数容量，并通过 expert parallelism 扩展到多设备。


\begin{importantbox}{什么是 MoE}
Mixture of Experts 把 dense MLP 替换成多个 experts，并由 router 为每个 token 选择少量 experts。总参数量可以很大，但每个 token 只激活少量参数，因此每 token FLOPs 可控。
\end{importantbox}

\begin{figure}[H]
\centering
\includegraphics[width=0.88\textwidth]{slide-013.jpg}
\caption{Slide 14: Mixture of experts}
\figsource{CS336 Spring 2026 Lecture 4 官方幻灯片。}
\end{figure}

\noindent\textbf{读图说明：}这张过渡页借当时业界对大模型采用 MoE 的讨论引出主题，但专有模型的具体结构没有公开证据，不能把“GPT4 (?)”当作事实引用。教学重点是区分两种规模：dense model 的总参数几乎都会参与每个 token 的计算；MoE 可以让总参数容量远大于单 token 激活参数，从而把“模型能存多少模式”和“每步花多少 FLOPs”部分解耦。


\begin{figure}[H]
\centering
\includegraphics[width=0.88\textwidth]{slide-014.jpg}
\caption{Slide 15: What’s a MoE?}
\figsource{CS336 Spring 2026 Lecture 4 官方幻灯片。}
\end{figure}

\noindent\textbf{读图说明：}Fedus 等人的示意图把 dense MLP 拆成 router 与多个 expert MLP。读图时先沿单个 token 的路径看：hidden state 进入 router，得到各 expert 的分数，只把 token 发送给 top-$k$ experts，再按 gate 权重合并输出。总参数量包含全部 experts，但 active parameters 只包含本次被选中的部分；这正是 MoE 提升容量而保持近似固定计算量的机制。


\begin{figure}[H]
\centering
\includegraphics[width=0.88\textwidth]{slide-015.jpg}
\caption{Slide 16: Why are MoEs getting popular?}
\figsource{CS336 Spring 2026 Lecture 4 官方幻灯片。}
\end{figure}

\noindent\textbf{展开说明：}Same FLOP, more param does better

\begin{knowledgebox}{读图：Slide 16 应该怎么看}
这页是机制或证据页。先看它比较的是 compute、参数量、routing、负载均衡还是通信；再看结论支持的是质量提升、训练速度、推理成本还是系统复杂度。MoE 和 attention alternatives 的图常把模型结构和系统代价混在一起，读图时要分清“算法机制”和“硬件执行成本”。
\end{knowledgebox}



\begin{figure}[H]
\centering
\includegraphics[width=0.88\textwidth]{slide-016.jpg}
\caption{Slide 17：MoE 在相同目标质量下可更快达到训练结果。}
\figsource{CS336 Spring 2026 Lecture 4 官方幻灯片。}
\end{figure}

\noindent\textbf{展开说明：}Faster to train MoEs。比较时必须固定 active FLOPs 和 token budget：MoE 通过每个 token 只激活少数 experts，让总参数容量远大于每步计算；曲线更快并不表示其总参数更少，也不保证 wall-clock 更快，因为 routing、all-to-all 和不均衡会吃掉理论收益。

\begin{knowledgebox}{讲义补充：源 Slide 17 应该怎么看}
这页是机制或证据页。先看它比较的是 compute、参数量、routing、负载均衡还是通信；再看结论支持的是质量提升、训练速度、推理成本还是系统复杂度。MoE 和 attention alternatives 的图常把模型结构和系统代价混在一起，读图时要分清“算法机制”和“硬件执行成本”。
\end{knowledgebox}



\begin{figure}[H]
\centering
\includegraphics[width=0.88\textwidth]{slide-017.jpg}
\caption{Slide 18：MoE 与 dense equivalent 的质量和成本竞争力。}
\figsource{CS336 Spring 2026 Lecture 4 官方幻灯片。}
\end{figure}

\noindent\textbf{展开说明：}Highly competitive vs dense equivalents。第一比较轴是相同 training compute 下的 loss，第二轴是相同 serving budget 下的 active parameters、显存和通信。MoE 的优势来自扩大可用知识容量，但部署时仍要存储所有 experts；因此“每 token FLOPs 小”与“模型容易服务”是两件不同的事。

\begin{importantbox}{MoE 证据的四本账}
阅读 Slides 16--24 时分别记录 total parameters、active parameters、training FLOPs 与 end-to-end time。Total parameters 决定容量和存储，active parameters 近似决定 token 计算，training FLOPs 用于公平 scaling 比较，end-to-end time 则包含 dispatch、all-to-all、capacity padding 与负载不均衡。只报告其中一项，都会把 conditional compute 的代价藏起来。
\end{importantbox}

\begin{knowledgebox}{讲义补充：源 Slide 18 应该怎么看}
这页是机制或证据页。先看它比较的是 compute、参数量、routing、负载均衡还是通信；再看结论支持的是质量提升、训练速度、推理成本还是系统复杂度。MoE 和 attention alternatives 的图常把模型结构和系统代价混在一起，读图时要分清“算法机制”和“硬件执行成本”。
\end{knowledgebox}


\begin{figure}[H]
\centering
\includegraphics[width=0.88\textwidth]{slide-018.jpg}
\caption{Slide 19: Why are MoEs getting popular?}
\figsource{CS336 Spring 2026 Lecture 4 官方幻灯片。}
\end{figure}

\noindent\textbf{展开说明：}Parallelizable to many devices

\begin{knowledgebox}{读图：Slide 19 应该怎么看}
这页是机制或证据页。先看它比较的是 compute、参数量、routing、负载均衡还是通信；再看结论支持的是质量提升、训练速度、推理成本还是系统复杂度。MoE 和 attention alternatives 的图常把模型结构和系统代价混在一起，读图时要分清“算法机制”和“硬件执行成本”。
\end{knowledgebox}


\begin{figure}[H]
\centering
\includegraphics[width=0.88\textwidth]{slide-019.jpg}
\caption{Slide 20: Some MoE results – from the west}
\figsource{CS336 Spring 2026 Lecture 4 官方幻灯片。}
\end{figure}

\noindent\textbf{展开说明：}MoEs are most of the highest-performance open models, and are quite quick

\begin{knowledgebox}{读图：Slide 20 应该怎么看}
这页是机制或证据页。先看它比较的是 compute、参数量、routing、负载均衡还是通信；再看结论支持的是质量提升、训练速度、推理成本还是系统复杂度。MoE 和 attention alternatives 的图常把模型结构和系统代价混在一起，读图时要分清“算法机制”和“硬件执行成本”。
\end{knowledgebox}


\begin{figure}[H]
\centering
\includegraphics[width=0.88\textwidth]{slide-020.jpg}
\caption{Slide 21: Earlier MoE results from Chinese groups – Qwen}
\figsource{CS336 Spring 2026 Lecture 4 官方幻灯片。}
\end{figure}

\noindent\textbf{展开说明：}Chinese LLM companies are also doing quite a bit of MoE work on the smaller end


\begin{figure}[H]
\centering
\includegraphics[width=0.88\textwidth]{slide-021.jpg}
\caption{Slide 22: Earlier MoE results from Chinese groups - DeepSeek}
\figsource{CS336 Spring 2026 Lecture 4 官方幻灯片。}
\end{figure}

\noindent\textbf{展开说明：}There’s also some good recent ablation work on MoEs showing they’re generally good


\begin{figure}[H]
\centering
\includegraphics[width=0.88\textwidth]{slide-022.jpg}
\caption{Slide 23: Recent MoE results – DeepSeek v3}
\figsource{CS336 Spring 2026 Lecture 4 官方幻灯片。}
\end{figure}

\noindent\textbf{展开说明：}Recent MoE results – DeepSeek v3

\begin{knowledgebox}{读图：Slide 23 应该怎么看}
这页是机制或证据页。先看它比较的是 compute、参数量、routing、负载均衡还是通信；再看结论支持的是质量提升、训练速度、推理成本还是系统复杂度。MoE 和 attention alternatives 的图常把模型结构和系统代价混在一起，读图时要分清“算法机制”和“硬件执行成本”。
\end{knowledgebox}


\begin{figure}[H]
\centering
\includegraphics[width=0.88\textwidth]{slide-023.jpg}
\caption{Slide 24: Why haven’t MoEs been more popular?}
\figsource{CS336 Spring 2026 Lecture 4 官方幻灯片。}
\end{figure}

\noindent\textbf{展开说明：}Infrastructure is complex / advantages on multi node

\begin{knowledgebox}{读图：Slide 24 应该怎么看}
这页是机制或证据页。先看它比较的是 compute、参数量、routing、负载均衡还是通信；再看结论支持的是质量提升、训练速度、推理成本还是系统复杂度。MoE 和 attention alternatives 的图常把模型结构和系统代价混在一起，读图时要分清“算法机制”和“硬件执行成本”。
\end{knowledgebox}

\subsection{本章小结}
本节的共同线索是 selective computation：通过结构选择让 token 不必总是访问所有历史或所有参数。但选择性越强，越需要额外机制保证信息流、负载均衡和训练稳定。

